% SPDX-License-Identifier: Apache-2.0
%
% Vreteno: an RV32I core in TxHDL, the first large design. Built by
% //docs:vreteno. One column, because it is mostly listings.
\documentclass[journal,onecolumn,9pt]{IEEEtran}

% One column: 80 columns fit at the body size, so nothing is reduced.
\newcommand{\listingsize}{\normalsize}
% SPDX-License-Identifier: Apache-2.0
%
% The house preamble, shared by every document in this directory. Loaded
% after \documentclass, before \title. Keeping it in one file is what
% stops two articles drifting apart in font, listing style or figure
% style.
% Computer Modern. IEEEtran selects Times (ptm) by default, so the face has
% to be chosen explicitly.
%
% Latin Modern is the Computer Modern variety used here, and the choice is
% forced rather than aesthetic. Under T1 encoding, plain `cmr` has no Type 1
% outlines unless cm-super is installed, and it is not in the pinned
% distribution, so pdflatex falls back to EC bitmaps and every glyph in the
% PDF becomes a Type 3 font. That was measured: the first build produced a
% document with 10 Type 3 fonts and no Type 1. Latin Modern is Computer
% Modern redrawn with a full T1 repertoire and real .pfb outlines, so the
% same design comes out scalable.
\usepackage[T1]{fontenc}
\usepackage[utf8]{inputenc}
\usepackage{lmodern}
% microtype is not in the pinned TeX distribution. emergencystretch alone
% keeps the two-column measure from producing overfull lines.
\setlength{\emergencystretch}{3em}

\usepackage{amsmath}
\usepackage{amssymb}
\usepackage{amsthm}
\usepackage{booktabs}
\usepackage{array}
% enumitem is not in the pinned TeX distribution; the list spacing below
% does what \setlist would have done.
\usepackage{float}
\usepackage{xcolor}
\usepackage{listings}
\usepackage{tikz}
\usetikzlibrary{arrows.meta,positioning,fit,backgrounds,calc}
\usepackage[hidelinks,bookmarks=true]{hyperref}

% Numbered, definition-style box for a rule the reader is meant to apply.
\theoremstyle{definition}
\newtheorem{rulebox}{Rule}
\newtheorem{example}{Example}

% Unnumbered asides.
\theoremstyle{remark}
\newtheorem*{tip}{Tip}
\newtheorem*{pitfall}{Pitfall}

% Tighter lists, in place of enumitem's \setlist.
\setlength{\topsep}{2pt}
\setlength{\itemsep}{1pt}
\setlength{\parsep}{0pt}

\newcommand{\code}[1]{\texttt{#1}}
\newcommand{\term}[1]{\emph{#1}}

% Syntax highlighting. Four roles, distinguishable in colour and still
% distinguishable in greyscale, because an IEEE article gets printed: the
% keyword is the only bold role, the comment is the only italic one, and the
% two remaining roles differ in lightness as well as in hue.
\definecolor{kwblue}{RGB}{20,60,150}
\definecolor{cmtgreen}{RGB}{90,120,90}
\definecolor{strred}{RGB}{150,50,40}
\definecolor{numgray}{gray}{0.45}
\definecolor{framegray}{gray}{0.70}
\definecolor{bgshade}{gray}{0.97}
% A darker fill for figure cells. The listing background has to stay very
% light to keep code readable; a figure cell has to be clearly shaded or the
% distinction it draws is invisible in print.
\definecolor{cellshade}{gray}{0.90}

% The listing size is the document's choice, because it depends on the
% column: a two-column article sets `\listingsize` to 6pt and keeps its
% listings within 70 columns, a one-column document keeps the body size
% and includes source files at 80. Lines are never broken; a line that does not fit
% overflows the frame, which is visible, rather than wrapping, which is
% not.
\providecommand{\listingsize}{\scriptsize}

% `'` is deliberately not a string delimiter: the language uses it for loop
% labels, as Rust does, so treating it as a quote would swallow the rest of
% the line.
\lstdefinestyle{txhdl}{
  basicstyle=\ttfamily\listingsize,
  keywordstyle=\bfseries\color{kwblue},
  commentstyle=\itshape\color{cmtgreen},
  stringstyle=\color{strred},
  numberstyle=\tiny\color{numgray},
  morekeywords={mod,use,pub,const,type,struct,enum,trait,impl,fn,async,let,mut,
    match,if,else,loop,while,for,in,break,continue,return,as,await,
    module,bus,channel,transaction,protocol,pipeline,flow,seq,config,
    port,stage,pipe,instance,connect,bind,tag,domain,blackbox,foreign,
    property,role,signal,handler,true,false,
    sequence,interface,modport,implements,package,import,var,wire,func,proc,
    case,default,next,fork,branch,join,state,fsm},
  morecomment=[l]{//},
  morestring=[b]",
  showstringspaces=false,
  columns=fullflexible,
  keepspaces=true,
  breaklines=false,
  % Framed, so a listing is bounded rather than floating in the column.
  frame=single,
  rulecolor=\color{framegray},
  framesep=3pt,
  backgroundcolor=\color{bgshade},
  xleftmargin=3pt,
  xrightmargin=3pt,
  aboveskip=6pt,
  belowskip=6pt,
  % Every listing is numbered and captioned, so it can be referred to by
  % number. `caption` without `float` numbers the listing and leaves it
  % where it was written, which is what a listing in running prose needs.
  captionpos=b,
  abovecaptionskip=2pt,
  belowcaptionskip=6pt,
}

% Figures drawn as text. Framed the same way, and with the highlighting
% switched off, because a box-and-arrow diagram is not code and half its
% words turning blue reads as an accident. Setting a style adds keys rather
% than clearing the ones already set, so the three roles are neutralised by
% hand rather than by leaving them out.
\lstdefinestyle{ascii}{
  basicstyle=\ttfamily\listingsize,
  keywordstyle=\ttfamily,
  commentstyle=\ttfamily,
  stringstyle=\ttfamily,
  columns=fullflexible,
  keepspaces=true,
  frame=single,
  rulecolor=\color{framegray},
  framesep=3pt,
  backgroundcolor=\color{bgshade},
  xleftmargin=3pt,
  xrightmargin=3pt,
  aboveskip=6pt,
  belowskip=6pt,
}
\lstset{style=txhdl}

% House TikZ styles. `shorten` lives on the arrow styles rather than on each
% arrow, so every arrow in the document gets the gap, including ones added
% later. TikZ clips a path at a node's border, so without it an arrowhead
% butts against the box with no gap at all. Keep `node distance` well above
% twice the shorten, or the arrow shrinks to nothing.
\tikzset{
  box/.style={draw,rounded corners=2pt,align=center,inner sep=4pt,
              font=\footnotesize},
  boxf/.style={box,fill=bgshade},
  lbl/.style={font=\scriptsize\itshape,align=center},
  fl/.style={-{Stealth[length=4pt]},shorten >=2.5pt,shorten <=2.5pt},
  fld/.style={fl,dashed},
}



\InputIfFileExists{buildstamp.tex}{}{}
\providecommand{\buildstamp}{Draft build (outside the Bazel flow).}

\title{Vreteno: An RV32I Core in TxHDL}

\author{Filip Filmar and Dragiša Janković%
\thanks{This document holds the Vreteno core, a RISC-V RV32I processor
written in TxHDL under \code{//cpu/vreteno} of the project repository,
included verbatim at build time with the reference model it is checked
against, the programs it runs, and the trace and waveform the build
produced by running it. The language is described in the embedding
article~\cite{embedding}, the runtime in its own
document~\cite{runtime}, the smaller designs in the
examples~\cite{examples}, and the cover~\cite{cover} lists all of them.}%
\thanks{\buildstamp}}

\markboth{Vreteno: An RV32I Core in TxHDL}{Filmar and Janković: Vreteno}

\begin{document}
\maketitle

\begin{abstract}
Vreteno is a 32-bit RISC-V core, the RV32I base integer set, written
in TxHDL: the first design large enough to ask whether the language
scales past its examples. This is its first form, a single-cycle
core in simulation, one process that fetches, decodes, executes and
writes back on every rising edge. It is checked in lockstep against
a reference model of the instruction set written as a plain program:
after every cycle the program counter, the thirty-one registers and
the halt must agree, and at the halt the data memory must. The
demonstration program and sixty-four random ones pass, and a
deliberately broken shift fails at its first use. Every listing here
is the file in the tree, and the trace and the figure are what the
build produced when it ran the core.
\end{abstract}

\begin{IEEEkeywords}
RISC-V, RV32I, processor, hardware description languages, Rust,
lockstep testing.
\end{IEEEkeywords}

{\footnotesize\tableofcontents}

\section{What Vreteno Is}
\label{sec:v-what}

\emph{Vreteno} is the Serbian word for a spindle: the part that turns
and on which the rest is wound. The core is the first design in TxHDL
with more than a handful of registers, and it exists to find out three
things: whether the language's constructs, a unit with registers and
memories, one process, waits, reads and deferred drives, are enough to
write a processor without stepping outside them; whether a design of
this size can be tested to a standard the small examples cannot reach;
and what the lowering must grow to before such a design becomes a
netlist.

The plan has three steps, and this document is the first. The core is
written in simulation form and checked against a reference model, so
the design is known correct before anything else is asked of it. Next
the core becomes a pipeline, two or three stages, and the lockstep
check is what keeps that refactoring honest. Last, the lowering grows
until the core is a netlist that simulates, under nvc, against the
same traces, which is the check the smaller
examples~\cite{examples} already pass.

\section{The Instruction Set as Bits}
\label{sec:v-isa}

Listing~\ref{lst:v-isa} is RV32I as encoders and a decoder. An
encoder per mnemonic builds the word, the six formats being six
functions of their fields; the decoder takes a word apart into its
mnemonic, its registers and its immediate, sign-extended. The core
does not use this decoder: it decodes on its own, in
Section~\ref{sec:v-core}, from the fields of the word with the
language's \code{slice}, \code{concat} and \code{sext}. So there are
two decoders, written differently, and the lockstep test compares
what they lead to; a defect in either shows up as a disagreement. The
disassembler is for traces a person reads, such as
Listing~\ref{lst:v-out}. A unit test checks each format's encoder
against the decoder.

\lstinputlisting[caption={\code{isa.rs}: the encoders, the decoder,
the disassembler.},label={lst:v-isa}]{cpu/vreteno/src/isa.rs}

\section{The Reference Model}
\label{sec:v-model}

Listing~\ref{lst:v-model} is the instruction set as a program. Its
state is the architectural state and nothing else: the program
counter, the registers, a data memory, and whether it has halted.
One \code{step} is one instruction, written against the decoder in
the plainest Rust that says what the manual says. It is the oracle:
the core is right when it agrees with this, cycle by cycle. The data
memory begins at \code{DATA\_BASE}, and a load or store outside it is
a fault, so a wild address in a random program is caught rather than
wrapped.

\lstinputlisting[caption={\code{model.rs}: the reference
model.},label={lst:v-model}]{cpu/vreteno/src/model.rs}

\section{The Core}
\label{sec:v-core}

Listing~\ref{lst:v-core} is the core. It is one unit whose state is a
program counter, a halt bit, and three memories: the register file,
the instruction memory, loaded from the program with
\code{Mem::with}, and the data memory. Its input is the reset and its
outputs are the halt, the word fetched, and a \code{Writeback} value,
the register written this cycle and what was written, which is what
the waveform decodes field by field.

The process is one \code{loop} with one wait, the rising edge, and
everything after the wait is one cycle: the fetch, a plain read of
the instruction memory at the program counter; the decode, the fields
of the word by \code{slice} and the five immediates assembled by
\code{concat} and \code{sext}, as the manual draws them; two plain
reads of the register file; the execute, an ALU shared by the
register and immediate forms and a branch condition, each a function
of its inputs; and the writeback, a deferred write to the register
file, a deferred write to the data memory, and a deferred drive of
the program counter. A store of a byte or a half reads the word,
merges the lanes and writes it back, which one read port and one
write port allow within a cycle. Register zero is never written, and
the reset takes the same path as an instruction: it is the first case
of the selection, and drives the program counter to zero.

This is the simulation form. The decode and the selection are Rust
\code{match}es over the fields, which the lowering does not reach
yet, and the ALU is a function that returns a value rather than a
\code{case!} that drives one. The state, the waits and the drives are
the language's, so the lockstep test and the waveform are real, and
the lowering, when it comes, is a change to this file and not a
rewrite: the \code{match}es become \code{case!} arms over enums, and
the functions become the operators they already call.

\lstinputlisting[caption={\code{core.rs}: the single-cycle
core.},label={lst:v-core}]{cpu/vreteno/src/core.rs}

\section{The Programs}
\label{sec:v-programs}

Listing~\ref{lst:v-program} holds the programs. \code{Asm} is the
smallest assembler that has labels: a branch to a label not yet
placed is emitted as a placeholder and patched when the label is,
which is all a forward branch needs. The demonstration program sums
one to ten in a loop, calls a function that doubles the sum and
returns through \code{jalr}, stores and loads every width so the sign
extension of \code{lb} and \code{lh} and the zero extension of
\code{lbu} and \code{lhu} are exercised, uses the upper immediates,
the arithmetic shift and both compares, and ends in \code{ebreak}
with 110 in \code{x10} and at the first data word.

The random generator writes straight-line programs of a given length
from a seed: every register operation, aligned stores and loads
within the first data words through \code{x2}, which it keeps at the
data base, and now and then a forward branch over one or two words,
so both outcomes of every branch condition occur. Nothing in it can
loop, so every program halts.

\lstinputlisting[caption={\code{program.rs}: the assembler, the
demonstration, the random
programs.},label={lst:v-program}]{cpu/vreteno/src/program.rs}

\section{The Lockstep Test}
\label{sec:v-lockstep}

Listing~\ref{lst:v-lockstep} is the test, \code{bazel test
//cpu/vreteno:lockstep}. It builds the core with the program, keeps
handles on its program counter, register file, data memory and halt
bit, which \code{Reg} and \code{Mem} give out as clones, holds the
reset for one cycle, and then runs the core and the model a cycle and
a step at a time. After each, the program counter, the thirty-one
registers and the halt must agree, and the failure names the cycle
and the instruction; at the halt every word of data memory must
agree too. The demonstration program passes with the values the
comments promise, and sixty-four random programs of two hundred
instructions each pass, which covers every instruction the generator
emits many times over with operands it did not choose.

A test that cannot fail proves nothing, so the negative control was
run: with the arithmetic shift in the ALU replaced by the logical
one, the demonstration fails at its \code{srai}, naming the register,
the cycle and the instruction, and the random programs fail within a
few cycles of their first \code{sra}. The control was then removed.

\lstinputlisting[caption={\code{lockstep.rs}: the core against the
model.},label={lst:v-lockstep}]{cpu/vreteno/tests/lockstep.rs}

\section{The Run}
\label{sec:v-run}

Listing~\ref{lst:v-main} runs the demonstration program and prints a
line per cycle, the program counter, the disassembly of the word
fetched and the register written; Listing~\ref{lst:v-out} is what it
printed when the build ran it. The loop is visible as the three
instructions that repeat ten times, the call as the jump to
\code{0x5c} and the return to \code{0x20}, and the loads as the sign
and zero extensions of the same byte and half. Figure~\ref{fig:v-wave}
is the opening of the same run, the reset and the first dozen
instructions, drawn from the waveform the run wrote; the whole run
is in the FST beside it, for a viewer.

\lstinputlisting[caption={\code{main.rs}: the run and the
waveform.},label={lst:v-main}]{cpu/vreteno/src/main.rs}

\begin{figure}[htbp]
\centering
\input{vreteno_timing.tex}
\caption{The opening cycles: reset, then an instruction per edge;
\code{rd} and \code{val} are the writeback, \code{rd} zero when
nothing was written.}
\label{fig:v-wave}
\end{figure}

\lstinputlisting[style=ascii,caption={What \code{vreteno} printed when
the build ran it.},label={lst:v-out}]{out_vreteno.txt}

\section{What Comes Next}
\label{sec:v-next}

The pipeline. A two-stage core, fetch and execute, keeps the same
outputs and the same test, with the instruction register the
boundary between the stages and the branch resolved a cycle late;
three stages add a writeback register and the forwarding it needs.
The lockstep test compares architectural state, so a pipeline that
retires one instruction per cycle after its fill passes the same
test with the model stepped when the core retires, which is the
change the test needs and the only one.

The lowering. The core uses \code{match} over fields where a lowered
unit would use \code{case!} over enums, functions where it would use
operators in arms, and memories, which the lowering does not emit
yet. Each is a change the examples document already has a
precedent for or a stated next step toward. When the core lowers,
the same trace that checks the model checks the netlist.

\begin{thebibliography}{10}

\bibitem{embedding}
F.~Filmar and D.~Janković, ``Deleting a language: Embedding TxHDL in
Rust,'' project repository \code{HDL/txhdl}, \code{//docs:embedding},
2026.

\bibitem{runtime}
F.~Filmar and D.~Janković, ``The TxHDL runtime,'' project repository
\code{HDL/txhdl}, \code{//docs:runtime}, 2026.

\bibitem{examples}
F.~Filmar and D.~Janković, ``TxHDL by example,'' project repository
\code{HDL/txhdl}, \code{//docs:examples}, 2026.

\bibitem{cover}
F.~Filmar and D.~Janković, ``TxHDL documents,'' project repository
\code{HDL/txhdl}, \code{//docs:cover}, 2026.

\end{thebibliography}

\end{document}
